% Integration copy; canonical source: contributions/arnav/sections/arnav_model.tex
\section{Problem and observation convention}
The project asks whether a mechanistic predator--prey model can explain recurring
hare--lynx fluctuations and support useful predictions beyond the observations used
for calibration. This section establishes the equations, their biological assumptions,
units and dimensionless form. It supplies the common mathematical foundation for the
team's stability analysis, numerical simulations and parameter estimation.

The completed historical iteration uses annual observations labelled 1900--1920.
The Stan case study supplies a public transcription of Hudson's Bay Company pelt
records in thousands of pelts, with a documented source chain
\cite{arnav_carpenter2018}. These measurements are trade records, not population
censuses. We interpret $H(t)$ and $P(t)$ as nonnegative hare and lynx abundance
\emph{indices}, respectively, in explicitly declared species-specific units.
The comparison between a continuous state and an annual pelt total is an observation
approximation: trapping effort, capture probability, seasonal aggregation and
unmodelled ecology can all affect the relationship.

Advay's frozen data split and the reference configuration use one index unit per
thousand pelts for each species \cite{arnav_advay2026,arnav_historical2026}.
The Stan compilation remains the source for this iteration. The alternative
tsibbledata compilation disagrees with it, a documented source limitation; neither
series establishes absolute animal abundance, and they are not combined.
Set $t=\mathrm{Year}-1900$ in years. The implemented evaluation design uses
1900--1915 for calibration and 1916--1920 for final scoring. Normalization and
parameter selection use training observations only. Forecasts have already been
frozen and scored; this section documents their observation convention without
changing it. It does not provide additional historical fits.

\subsection{Explicit observation model}
For year $\mathrm{Year}_i$ and $t_i=\mathrm{Year}_i-1900$, write
\begin{equation}\label{eq:arnav_observation}
Y_{H,i}=H(t_i)+\varepsilon_{H,i},\qquad
Y_{P,i}=P(t_i)+\varepsilon_{P,i}.
\end{equation}
Here $Y$ is the recorded annual pelt count in thousands; $H,P$ are the model's
noiseless pelt indices; and $\varepsilon$ is observation discrepancy in the same
species-specific units. Under the working proportional-index assumption O1,
\begin{equation}\label{eq:arnav_O1}
H(t)=c_H N_H(t),\qquad P(t)=c_P N_P(t),\qquad c_H,c_P>0,
\end{equation}
where $N_H,N_P$ denote latent animal abundance and the unknown observation factors
are constant. We fit the effective equations directly in pelt units:
\emph{no observation-scale factor is fitted}. This does not set $c_H=c_P=1$ or
identify actual animal counts. Constant factors are absorbed into the effective
coefficients through the transformation derived in Eq.~\eqref{eq:arnav_unitmap}.

Each annual total is approximated by a point observation of the state, as in the
implemented convention \cite{arnav_advay2026,arnav_historical2026}. Thus $H,P$ have
units of thousands of pelts, and $\dot H,\dot P$ have units of thousands of pelts
per year. No seasonal observation integral is fitted. A latent-population model
with estimated observation factors, or an integral over the trapping season,
would represent additional mechanisms but require further assumptions and
unknowns. O1 and the point-sample approximation keep the present comparison
tractable, at the cost of limited biological interpretation.

The main observation loss compares additive residuals after division by the
training maxima $s_H=77.4$ and $s_P=59.4$:
\begin{equation}\label{eq:arnav_obsresidual}
R_{H,i}=\frac{H(t_i)-Y_{H,i}}{s_H},\qquad
R_{P,i}=\frac{P(t_i)-Y_{P,i}}{s_P}.
\end{equation}
For the PINN data term the state in this expression is the network approximation;
for least squares it is the integrated ODE state. Both methods use the same
species scales. These weights are a declared comparison convention, not measured
noise standard deviations. Additive discrepancy notation alone establishes
neither Gaussian noise nor independence of observations. The reported log-error
sensitivity instead emphasizes relative discrepancies. A separate sensitivity
fits the 1900 initial states within $[0.5,1.5]$ times their observed values
\cite{arnav_historical2026}; the two checks change different assumptions and
remain separate outcomes.

\section{From Lotka--Volterra to the chosen model}
\subsection{The reference model and its limitation}
For a first comparison, let
\begin{equation}\label{eq:arnav_lv}
\dot H=rH-cHP,\qquad \dot P=ecHP-mP,
\end{equation}
where a dot means differentiation with respect to years. The prey grows at rate
$rH$ without predators, encounters remove prey at rate $cHP$, conversion of consumed
prey contributes $ecHP$ to predator growth, and predator mortality removes $mP$.
Here $c>0$ is the encounter coefficient; it is not the saturation parameter $a$
introduced below. The other parameters are positive.

The positive equilibrium is $H_*=m/(ec)$, $P_*=r/c$. To see why the cycles are
neutral, rather than simply asserting it, define for positive $H,P$
\begin{equation}
V=ec\left[H-H_*-H_*\log(H/H_*)\right]
 +c\left[P-P_*-P_*\log(P/P_*)\right].
\end{equation}
The logarithms have dimensionless arguments. Differentiation gives
\begin{align}
\dot V&=(ec-m/H)\dot H+(c-r/P)\dot P\\
&=(ecH-m)(r-cP)+(cP-r)(ecH-m)=0.
\end{align}
Each bracket in $V$ has a strict minimum at its equilibrium value and diverges at
zero or infinity. Thus nonzero level sets are closed curves in the positive
quadrant, surrounding the equilibrium. The vector field has no other positive
equilibrium on a level set, so trajectories follow periodic orbits. A perturbation
generally moves the state to a different orbit; it does not relax back to one
selected cycle. The equilibrium is a centre, not an attractor. Also, predator-free
prey grows exponentially without a resource limit. These are reasons to refine the
model; they do not mean that Lotka--Volterra cannot oscillate.

\subsection{Resource limitation and saturating predation}
Let the prey's per-capita net growth decrease linearly from $r$ at low abundance
to zero at $H=K$. Multiplication by $H$ gives logistic growth,
\begin{equation}
G(H)=rH(1-H/K).
\end{equation}
The carrying capacity $K$ describes the prey-only resource limit under fixed
environmental conditions; it is not an absolute ceiling on every trajectory.

Let $g(H)$ be consumption per predator. A simple search-and-handling argument
assumes an encounter-limited rate $qH$ while searching and a maximum processing
rate $a$. If the fractions of time searching and handling are $g/(qH)$ and $g/a$,
their sum is one. For $H>0$,
\begin{equation}
\frac{g}{qH}+\frac{g}{a}=1
\quad\Longrightarrow\quad
g(H)=\frac{aqH}{a+qH}=\frac{aH}{b+H},\qquad b=\frac{a}{q}.
\end{equation}
Extend continuously by $g(0)=0$. This is the Holling type II response
\cite{arnav_holling1959}. It assumes an effectively constant prey level over the
short feeding interval and no predator interference. In index units it is an
effective response law, not a direct measurement of handling time. The parameters
have transparent interpretations:
\begin{equation}
g(b)=a/2,\qquad \lim_{H\to\infty}g(H)=a,\qquad
g(H)\simeq(a/b)H\quad\text{when }H/b\ll1.
\end{equation}
The auxiliary $q=a/b$ has units inverse lynx index per year. It is not an additional
parameter to estimate. The linear response would be simpler but cannot represent
saturation; a type III response would introduce a different low-prey mechanism
that the present project does not establish.

\subsection{Dimensional Rosenzweig--MacArthur equations}
Combining logistic growth, total prey consumption $g(H)P$, conversion and mortality
gives the project's Rosenzweig--MacArthur model
\cite{arnav_rosenzweig1963}:
\begin{align}
\frac{dH}{dt}&=rH\left(1-\frac{H}{K}\right)-\frac{aHP}{b+H},\label{eq:arnav_H}\\
\frac{dP}{dt}&=\frac{eaHP}{b+H}-mP.\label{eq:arnav_P}
\end{align}
These are exactly the dimensional equations in the shared repository
\cite{arnav_teamrepo2026}. Logistic
growth and saturation allow stability to depend on parameters. Whether a specified
case approaches an equilibrium or exhibits a stable cycle requires the subsequent
analysis and numerical evidence; it does not follow merely from naming the model.

\section{Units and assumptions}
Let $U_H$ and $U_P$ denote the chosen hare and lynx index units. Keeping these labels
distinct prevents an unjustified interpretation of conversion efficiency.
\begin{center}
\small
\begin{tabularx}{\linewidth}{@{}lXl@{}}
\toprule
Symbol & Meaning & Units\\\midrule
$t$ & Elapsed time & year\\
$H$, $P$ & Hare and lynx states & $U_H$, $U_P$\\
$r$ & Low-density hare net growth rate & year$^{-1}$\\
$K$ & Effective prey-only carrying-capacity index & $U_H$\\
$a$ & Effective saturated consumption coefficient & $U_H U_P^{-1}$ year$^{-1}$\\
$b$ & Prey index at half saturation & $U_H$\\
$e$ & Effective predator conversion coefficient & $U_P U_H^{-1}$\\
$m$ & Predator mortality rate & year$^{-1}$\\
$H_0$, $P_0$ & States at the initial time & $U_H$, $U_P$\\
\bottomrule
\end{tabularx}
\end{center}
All six parameters are strictly positive. For example,
$[aHP/(b+H)]=U_H$/year and $[eaHP/(b+H)]=U_P$/year, matching the two left-hand
sides. In separately scaled abundance indices, $e$ is an effective conversion
factor. It need not be a dimensionless biomass efficiency bounded above by one.
For this historical iteration $U_H$ means thousands of hare pelts and $U_P$
thousands of lynx pelts. The rates $r,m$ retain their growth/mortality interpretation
only conditional on O1 and an adequate dynamical model; $K,b,a,e$ are effective
pelt-unit coefficients. Scale invariance does not establish accurate recovery of
any biological parameter.

The simplifying assumptions and their consequences are:
\begin{enumerate}
\item A closed two-species system has no immigration, emigration, alternative prey
or other predators. This isolates the proposed feedback but omits a food web.
\item Continuous, spatially mixed states omit age structure, territories and local
differences. The approximation is unsuitable for individual-level extinction claims.
\item Constant parameters omit seasons, weather and changing resource availability.
Fitted rates summarize the chosen interval rather than universal ecological constants.
\item Resource limitation is represented by logistic prey growth. Plants are not
modelled as a third dynamical state.
\item Predation follows the type II response with no interference between predators.
Predator growth depends on prey; without prey, $P$ decays at rate $m$.
\item Pelt totals are treated as noisy abundance proxies. Constant observation
scales are a working assumption. Unknown or changing capture rates can confound
parameter inference. Harvest is not modelled as an additional removal process.
\end{enumerate}

\subsection{Basic consistency checks}
The equations have the form $\dot H=H A(H,P)$, $\dot P=P B(H)$ with continuous
coefficients for $H\geq0$. Along a finite solution interval,
\begin{equation}
H(t)=H_0\exp\!\left(\int_0^t A(H(s),P(s))\,ds\right),\quad
P(t)=P_0\exp\!\left(\int_0^t B(H(s))\,ds\right).
\end{equation}
Positive initial states remain positive, and either zero boundary is invariant by
uniqueness. A numerical method can still introduce negative approximations, so its
outputs must be checked. Also $\dot H\leq rH(1-H/K)$, giving
$H(t)\leq\max\{H_0,K\}$ for $t\geq0$. Since $\dot P\leq eaP$, predator abundance
cannot blow up in finite time. These checks establish basic consistency, not
ecological accuracy or long-term stability.

\section{Non dimensionalisation}
\subsection{Choosing the scales}
For a fixed parameter vector define
\begin{equation}\label{eq:arnav_scales}
s=rt,\qquad x=H/K,\qquad y=\frac{aP}{rK}.
\end{equation}
The corresponding reference time and states are $1/r$, $K$ and $rK/a$.
Their units are year, $U_H$ and $U_P$. The scales therefore make all three new
variables dimensionless. Choosing $K$ for prey makes the prey-only capacity one;
choosing $rK/a$ for predators makes the prey-loss coefficient one. This is the
same notation used in Joel's existing dimensionless analysis
\cite{arnav_joel2026}.

Define the three dimensionless groups
\begin{equation}\label{eq:arnav_groups}
\beta=b/K,\qquad \alpha=ea/r,\qquad \delta=m/r.
\end{equation}
They compare half saturation with carrying capacity, maximum predator growth with
prey growth, and predator mortality with prey growth, respectively. This scaling
also means a sweep of $K$ changes both $\beta$ and the state scales. If the
dimensional initial states are held fixed in that sweep, their scaled values must
be recomputed for each $K$.

\subsection{Substitution and chain rule}
Because $H=Kx$, $P=(rK/a)y$ and $ds/dt=r$,
\begin{equation}
\frac{dH}{dt}=rK\frac{dx}{ds},\qquad
\frac{dP}{dt}=\frac{r^2K}{a}\frac{dy}{ds}.
\end{equation}
Substitution into the prey equation gives
\begin{equation}
rK\frac{dx}{ds}=rKx(1-x)-\frac{rKxy}{b/K+x}.
\end{equation}
Substitution into the predator equation gives
\begin{equation}
\frac{r^2K}{a}\frac{dy}{ds}
=\frac{erKxy}{b/K+x}-\frac{mrK}{a}y.
\end{equation}
Dividing by the respective factors yields
\begin{align}
\frac{dx}{ds}&=x(1-x)-\frac{xy}{\beta+x},\label{eq:arnav_x}\\
\frac{dy}{ds}&=\alpha\frac{xy}{\beta+x}-\delta y.\label{eq:arnav_y}
\end{align}
The two systems describe the same trajectories after the changes of units. Reducing
six dimensional parameters to three groups does not prove that all six can be
identified from data: physical scales and initial conditions also carry information.

\subsection{Initial conditions and conversion back}
At year 1900, $s=0$, $x_0=H_0/K$ and $y_0=aP_0/(rK)$. Recover the dimensional
variables using
\begin{equation}
\mathrm{Year}=1900+s/r,\qquad H=Kx,\qquad P=(rK/a)y.
\end{equation}
If $r,K,a$ are separately specified, recover
$b=\beta K$, $e=\alpha r/a$ and $m=\delta r$. The three dimensionless groups alone
cannot reconstruct all six dimensional parameters.

For an arithmetic check, use the repository's synthetic values
\cite{arnav_teamrepo2026}:
\[
(r,K,a,b,e,m)=(1,60,2,20,0.6,0.5),\qquad (H_0,P_0)=(35,9).
\]
These give
\begin{equation}
(\beta,\alpha,\delta)=(1/3,1.2,0.5),\qquad (x_0,y_0)=(7/12,0.3).
\end{equation}
At this initial state, $(\dot H,\dot P)=(413/132,261/110)$ in their respective
index units per year. The dimensionless derivatives are
$(dx/ds,dy/ds)=(413/7920,87/1100)$; multiplying by $rK=60$ and
$r^2K/a=30$ recovers the dimensional values. These are synthetic checks, not fitted
historical estimates.

An equally valid alternative is $X=H/b=x/\beta$ and $Y=aP/(rb)=y/\beta$,
with the same time $s$. Writing $\kappa=K/b=1/\beta$ gives
$X_s=X(1-X/\kappa)-XY/(1+X)$ and
$Y_s=\alpha XY/(1+X)-\delta Y$. This keeps state scales fixed during a $K$
sweep when other parameters remain fixed. We use the $K$-based form as the main
convention to match the existing team analysis, while retaining the alternative
as an explanation of the scaling tradeoff.

\section{Compatibility with data scaling and the PINN}
\subsection{Keep mathematical and numerical coordinates explicit}
The historical PINN \cite{arnav_historical2026} uses fixed training maxima $s_H,s_P$, normalized outputs
$u=H/s_H$, $v=P/s_P$, training duration $T=15$ years, and internal time input
$\xi=2t/T-1$. It differentiates with respect to elapsed years $t$.
These numerical choices can coexist with the dimensionless theory coordinates.
In elapsed time,
\begin{align}
\frac{du}{dt}&=ru(1-s_Hu/K)-\frac{a s_Puv}{b+s_Hu},\\
\frac{dv}{dt}&=\frac{ea s_Huv}{b+s_Hu}-mv.
\end{align}
The bridge to the theoretical variables is
\begin{equation}
x=(s_H/K)u,\qquad y=\frac{a s_P}{rK}v,\qquad
s=\frac{rT}{2}(\xi+1).
\end{equation}
If differentiation is instead performed with respect to $\xi$, the chain rule
requires $du/dt=(2/T)du/d\xi$, and likewise for $v$. Thus a residual formed in
years is $R_u=(2/T)u_\xi-f_H/s_H$. If one forms a residual in $\xi$ coordinates,
$\widetilde R_u=u_\xi-(T/2)f_H/s_H=(T/2)R_u$. With the same averaging convention,
its squared penalty is $(T/2)^2$ larger; the loss coefficient must be multiplied
by $(2/T)^2$ to preserve the same objective. This is a unit change, not evidence
that one trained model is better.

Because all six coefficients, including $r$ and $a$, are inferred in this iteration, directly using $rt$
and $aP/(rK)$ as fixed training inputs or targets would introduce parameter-dependent
coordinates. Retaining his fixed training normalization avoids that complication.
The dimensionless form is for derivation, interpretation and consistent
mapping, and does not require rewriting the PINN.

\subsection{Why measurement units change parameter meanings}
For any constant positive changes $\widetilde H=c_HH$,
$\widetilde P=c_PP$, keeping time in years preserves the model form if
\begin{equation}\label{eq:arnav_unitmap}
\widetilde r=r,\quad \widetilde m=m,\quad
\widetilde K=c_HK,\quad \widetilde b=c_Hb,\quad
\widetilde a=\frac{c_H}{c_P}a,\quad
\widetilde e=\frac{c_P}{c_H}e.
\end{equation}
This follows by substituting into both derivatives, including the denominator.
It leaves $\beta,\alpha,\delta,x,y$ unchanged. For example, converting both species
from counts to thousands uses $c_H=c_P=10^{-3}$: $K$ and $b$ divide by 1000,
while $a$ and $e$ keep their numerical values. Unequal species scales change
$a$ and $e$ too. Unknown pelt-to-population multipliers have the same algebraic
effect, but their unknown values prevent conversion to absolute biological rates.
Time-varying multipliers would add derivative terms and are outside this model.

The same transformation leaves Joel's stability ratio $\sigma=K/K_H$ unchanged,
where $K_H=b(ea+m)/(ea-m)$ and positive coexistence is applicable
\cite{arnav_joel2026}. Thus the dimensionless groups have a model-based
interpretation independent of constant observation scaling. This does not confirm
that the real ecosystem follows the model or has a stable limit cycle.

With the observed predator scale fixed, an arbitrary rescaling of $P$ cannot
preserve the same observations. Moreover, $a$ occurs separately in the prey
equation while $ea$ occurs in the predator equation. Therefore the scale ambiguity
alone does not prove that effective $a$ and $e$ are inseparable. It also does not
prove that all six coefficients are well determined by 16 annual pairs.

\subsection{Reference settings for the completed iteration}
The reference is \texttt{configs/historical.json}, rather than the earlier
synthetic configuration \cite{arnav_historical2026}. State order is $(H,P)$ and
parameter names are $(r,K,a,b,e,m)$, matching the shared right-hand side.
All six coefficients are fitted. Their independent positive bounds are
\[
r,m\in[0.05,3],\quad K\in[20,500],\quad a\in[0.05,10],\quad
b\in[0.5,300],\quad e\in[0.02,5].
\]
The units are those in the table above. These are exploratory numerical bounds,
not measured ecological ranges. The completed iteration imposes neither $e\leq1$
nor a coexistence constraint. This describes the recorded computation, rather
than a new assertion of identifiability or a proposed change to the inverse problem.

Historical fitting and ODE forecasts use \texttt{model.solve}: RK45 with relative
tolerance $10^{-8}$, absolute tolerance $10^{-10}$ and maximum step $0.1$ years.
Rishi's tighter settings for illustrative long-run simulations are a separate
calculation, not the settings of this historical comparison
\cite{arnav_historical2026,arnav_rishi2026}.

In the main analysis the ODE starts at the observed 1900 index state $(30,4)$.
The PINN enforces its initial state with a soft penalty and reports the discrepancy;
we do not describe this as an exact network constraint. In the fitted-initial-state
sensitivity, each method forecasts from its fitted state at 1900. Every forecast
integrates from 1900 through 1920 with frozen parameters; none restarts from a
test-year observation. The historical handoff records both sensitivity variants
and preserves them alongside the main result \cite{arnav_historical2026}.

\section{Validation responsibilities and limits}
The derivations supply a model to test, not a claim that a published model is novel
or that a PINN earns bonus marks automatically. The independent contribution here
is the explicit reasoning, unit audit, coordinate transformations and interpretation
under the chosen observation convention. The accompanying notebook checks algebraic
equivalence at multiple exact rational inputs. These checks do not replace empirical
validation or establish parameter identifiability.

Joel's contribution supplies equilibrium, stability and Hopf analysis
\cite{arnav_joel2026}; Rishi's supplies numerical checks and parameter-regime
illustrations \cite{arnav_rishi2026}. Advay's data records and Yash's historical
fits, matching least-squares benchmark and evaluation are available
\cite{arnav_advay2026,arnav_historical2026}. Their empirical results belong in the
corresponding report sections. Comparisons use the same training years,
parameter bounds, initial-state convention and solver settings, with separate
hare and lynx MAE and RMSE in thousands of pelts. The PINN's soft initial penalty
and its network--ODE discrepancy must remain visible qualifications.

Five held-out years offer limited evidence. Synthetic recovery generated by the
same equations checks the inverse procedure under favourable assumptions; it is
not independent evidence that those equations describe real ecology. A small
training error does not show unique biological parameters. Any substantive change
after the completed evaluation must be documented as a later revision, preserving
the existing results. Report claims must remain restricted to the mathematical
and empirical evidence actually obtained.
